% U5 WS3 -- mechanisms, collision model, energy profiles. Blocks 4-5.
\documentclass{apchem-worksheet}

\apcunit{5}
\apcunittitle{Kinetics}
\apcdoctitle{Worksheet 3 \textbullet\ Mechanisms and Energy Profiles}
\apcsubtitle{CED 5.4--5.8 \textbullet\ molecularity gives the order only for an elementary step}

\begin{document}
\wsheader[Elementary step: coefficients \emph{are} orders \textbullet\
$k = Ae^{-E_a/RT}$ \textbullet\
$\ln(k_2/k_1) = -\tfrac{E_a}{R}\left(\tfrac{1}{T_2}-\tfrac{1}{T_1}\right)$
\textbullet\ $R = \SI{8.314}{\joule\per\mole\per\kelvin}$ \textbullet\
$E_a(\text{fwd}) - E_a(\text{rev}) = \Delta H$.]

\problem[]{\ced{5.4} Write the rate law for each \emph{elementary} step:}
\begin{parts}
  \item $\ce{A -> products}$ \answerblank[3.4cm]{rate $= k[\ce{A}]$}
  \item $\ce{A + B -> products}$
        \answerblank[3.4cm]{rate $= k[\ce{A}][\ce{B}]$}
  \item $\ce{2A -> products}$ \answerblank[3.4cm]{rate $= k[\ce{A}]^2$}
  \item $\ce{2NO + O2 -> 2NO2}$ as a single step
        \answerblank[4.4cm]{rate $= k[\ce{NO}]^2[\ce{O2}]$}
  \item Why is (d) an unlikely elementary step?
        \answerlines[2]{It would be termolecular --- three particles
        colliding simultaneously with the correct orientation --- which is
        so improbable that such steps are rarely proposed.}
\end{parts}

\problem[]{\ced{5.4} Explain precisely why coefficients give orders for an
elementary step but not for an overall equation.
\answerlines[4]{An elementary step \emph{is} the molecular event: its
coefficients state how many particles must actually collide, and the
probability of that collision scales with each concentration to that power.
An overall equation only summarizes the net stoichiometry of a sequence of
steps, and the observed rate is set by the slowest of them --- which may
involve entirely different species in different numbers.}}

\problem[]{\ced{5.7}--\ced{5.8} Overall: $\ce{2NO2 + F2 -> 2NO2F}$,
experimental rate $= k[\ce{NO2}][\ce{F2}]$. Proposed mechanism:
step 1 (slow) $\ce{NO2 + F2 -> NO2F + F}$; step 2 (fast)
$\ce{NO2 + F -> NO2F}$.}
\begin{parts}
  \item Show that the steps sum to the overall equation.
        \answerlines[2]{Adding: $\ce{2NO2 + F2 + F -> 2NO2F + F}$. The F
        atom appears on both sides and cancels, leaving
        $\ce{2NO2 + F2 -> 2NO2F}$.}
  \item Give the predicted rate law and say why.
        \answerlines[2]{Step 1 is rate-determining and elementary, so the
        rate law is that of step 1: rate $= k[\ce{NO2}][\ce{F2}]$.}
  \item Does the mechanism pass both tests?
        \answerblank[3cm]{yes}
  \item Identify the intermediate.
        \answerblank[2.6cm]{\ce{F}}
  \item Does passing both tests prove the mechanism is correct? Explain.
        \answerlines[3]{No. It shows the mechanism is \emph{consistent} with
        the evidence. Other mechanisms can predict the same rate law and sum
        to the same equation, so kinetics can rule mechanisms out but never
        single one in.}
\end{parts}

\problem[]{\ced{5.8} For the mechanism $\ce{A2 -> 2A}$ (slow), then
$\ce{A + B -> AB}$ (fast):}
\begin{parts}
  \item Overall equation:
        \answerblank[5cm]{$\ce{A2 + 2B -> 2AB}$}
  \item Predicted rate law: \answerblank[3.4cm]{rate $= k[\ce{A2}]$}
  \item B is a reactant but absent from the rate law. Is that a problem?
        \answerlines[2]{No. B appears only in the fast step, which is not
        rate-determining, so its concentration does not affect the overall
        rate. Zero order in a reactant is a perfectly normal result.}
\end{parts}

\problem[]{\ced{5.5} The collision model.}
\begin{parts}
  \item List the three requirements for a reactive collision.
        \answerlines[2]{The particles must collide; the collision must carry
        at least the activation energy; and the particles must be correctly
        oriented.}
  \item Why is the fraction of successful collisions so small?
        \answerlines[2]{Most collisions are either too low in energy to
        cross the barrier or occur at an unproductive orientation, so only a
        small tail of the distribution reacts.}
  \item Raising the temperature increases the rate far more than it
        increases collision frequency. Explain.
        \answerlines[3]{Higher temperature broadens the distribution of
        kinetic energies. Because the reactive fraction lies in the
        high-energy tail, a modest broadening moves a disproportionately
        large number of particles past $E_a$, which dominates the small
        increase in how often particles meet.}
\end{parts}

\problem[]{\ced{5.5} Arrhenius calculations.}
\begin{parts}
  \item A reaction with $E_a = \SI{50.0}{\kilo\joule\per\mole}$ is warmed
        from \SI{300}{\kelvin} to \SI{310}{\kelvin}. By what factor does
        $k$ change? \workspace[2.8cm]
        \keyonly{\begin{apcanswer}$\ln(k_2/k_1) = -(50000/8.314)
        (1/310 - 1/300) = 0.647$;
        $k_2/k_1 = e^{0.647} = \mathbf{1.9}$ --- roughly
        double\end{apcanswer}}
  \item Repeat for $E_a = \SI{100.0}{\kilo\joule\per\mole}$.
        \answerblank[3cm]{$\times 3.6$}
  \item What general conclusion follows?
        \answerlines[2]{The larger the activation energy, the more
        sensitive the rate is to temperature, because a bigger barrier means
        a smaller and more sharply changing reactive fraction.}
  \item A reaction has $k = \SI{2.0e-3}{\per\second}$ at \SI{300}{\kelvin}
        and \SI{8.0e-3}{\per\second} at \SI{320}{\kelvin}. Find $E_a$.
        \workspace[2.6cm]
        \keyonly{\begin{apcanswer}$\ln(8.0/2.0) = 1.386 =
        -(E_a/8.314)(1/320 - 1/300)$;
        $E_a = \mathbf{5.5\times10^{4}}~\si{\joule\per\mole} =
        \SI{55}{\kilo\joule\per\mole}$\end{apcanswer}}
\end{parts}

\problem[]{\ced{5.6} Energy profiles.}
\begin{parts}
  \item $E_a(\text{fwd}) = \SI{80}{\kilo\joule}$,
        $E_a(\text{rev}) = \SI{120}{\kilo\joule}$. Find $\Delta H$ and
        classify. \answerblank[5cm]{$-40$ kJ; exothermic}
  \item $E_a(\text{fwd}) = \SI{150}{\kilo\joule}$,
        $\Delta H = \SI{+60}{\kilo\joule}$. Find $E_a(\text{rev})$.
        \answerblank[3cm]{\SI{90}{\kilo\joule}}
  \item On a profile, $E_a$ is measured from what to what?
        \answerblank[7.5cm]{reactants up to the transition state}
  \item Can a strongly exothermic reaction be very slow? Explain.
        \answerlines[3]{Yes. $\Delta H$ and $E_a$ are independent: the
        products may lie far below the reactants while the barrier between
        them is still very high. Diamond converting to graphite is the
        standard case --- favorable, and effectively never observed.}
\end{parts}

\begin{spiralstrip}{Unit 5 \textbullet\ blocks 1--3}
  \item First-order half-life: \answerblank[2cm]{$0.693/k$}
  \item Which plot is linear for second order?
        \answerblank[3cm]{$1/[\ce{A}]$ vs $t$}
  \item Constant half-life implies: \answerblank[2.6cm]{first order}
  \item Orders come from: \answerblank[2.4cm]{experiment}
  \item Zero-order half-life: \answerblank[2.4cm]{$[\ce{A}]_0/2k$}
\end{spiralstrip}

\end{document}
